% Problem record op_6f6de1a285416bb2. This TeX file is derived from
% database/problems_json/op_6f6de1a285416bb2.json by scripts/sync-tex.mjs; edit the JSON record
% and rerun the script rather than editing this file.
\section{Optimal success probability for Gaussian processing of non-Gaussian entanglement}
\label{sec:op_6f6de1a285416bb2}

\paragraph{Problem.}
What is the maximal success probability for preparing an approximate two-mode squeezed vacuum from finitely many copies of a non-Gaussian state?

Let $\rho$ be a non-Gaussian density operator of one bosonic mode at Alice and one at Bob. Fix an integer $m\geq1$, squeezing $r>0$, and error tolerance $0<\varepsilon<1$. The target is the two-mode squeezed vacuum

\begin{equation}
\psi_r:=|\psi_r\rangle\langle\psi_r|,\qquad
|\psi_r\rangle:=\sqrt{1-\tanh^2r}\sum_{j=0}^{\infty}(\tanh r)^j|j,j\rangle,
\label{eq:ser18-statement-1}
\end{equation}

Equation~\eqref{eq:ser18-statement-1} uses the product photon-number basis.

Preparation from $m$ copies of $\rho$ uses local Gaussian operations and classical communication. Allowed success maps $\Lambda$ are trace-nonincreasing completely positive maps. Each implementation is a finite sequence using local Gaussian ancillary states, local Gaussian unitaries, Gaussian measurements or vacuum-projection success branches, discarding, and classical feedforward. Gaussian measurement records may be accepted on a specified measurable set; their probabilities are integrated over that set. A nonvacuum outcome of a vacuum projection terminates that run in failure; its output cannot be reused. The input copies are the only non-Gaussian resource.

For each allowed map, write $p_\Lambda:=\operatorname{Tr}\Lambda(\rho^{\otimes m})$. When $p_\Lambda>0$, let $\tau_\Lambda:=\Lambda(\rho^{\otimes m})/p_\Lambda$. Define the maximal heralding probability by

\begin{equation}
P_G^{(m)}(\rho,r,\varepsilon):=
\sup_{\Lambda}\left\{p_\Lambda:p_\Lambda>0,\quad
\tfrac12\|\tau_\Lambda-\psi_r\|_1\leq\varepsilon\right\}.
\label{eq:ser18-statement-2}
\end{equation}

Determine Eq.~\eqref{eq:ser18-statement-2}, with $\sup\varnothing:=0$. Here $\|X\|_1:=\operatorname{Tr}\sqrt{X^\dagger X}$.

\subsection*{Status}

Unsolved

\subsection*{Source}

This optimization is a precise formulation of the probability-versus-accuracy question suggested by Browne et al., Sections II--III and Figures 2--3. The source supplies particular protocols, not this universal optimum \sourcecite{ref:ser18-1}{BESP03}.

\subsection*{Progress}

\begin{itemize}
  \item For the input $|\varphi_\lambda\rangle:=(|0,0\rangle+\lambda|1,1\rangle)/\sqrt{1+\lambda^2}$, $0<\lambda<1$, a balanced beam-splitter tree with $m=2^k$ input copies, $k\geq0$, and vacuum projections on all discarded ports produces coefficients

\begin{equation}
\begin{aligned}
b_{m,j}&:=\lambda^j\prod_{\ell=0}^{j-1}(1-\ell/m),
&Z_m&:=\sum_{j=0}^{m}b_{m,j}^2,\\
|\varphi_{\lambda,m}\rangle&:=Z_m^{-1/2}
\sum_{j=0}^{m}b_{m,j}|j,j\rangle,
&p_m&:=\frac{Z_m}{(1+\lambda^2)^m}.
\end{aligned}
\label{eq:ser18-progress-1-1}
\end{equation}

Equation~\eqref{eq:ser18-progress-1-1} follows by expanding the published optical iteration; the empty product equals one. \sourcecite{ref:ser18-1}{BESP03}

  \item Consequently, with $r=\operatorname{arctanh}\lambda$ and $\varepsilon_m:=\sqrt{1-|\langle\psi_r|\varphi_{\lambda,m}\rangle|^2}$, the protocol gives

\begin{equation}
\begin{gathered}
P_G^{(m)}(|\varphi_\lambda\rangle\langle\varphi_\lambda|,r,\varepsilon_m)\geq p_m,
\\ \varepsilon_m\longrightarrow0,\qquad
p_m\sim\frac{(1+\lambda^2)^{-m}}{1-\lambda^2}.
\end{gathered}
\label{eq:ser18-progress-2-1}
\end{equation}

The limits in Eq.~\eqref{eq:ser18-progress-2-1} follow directly from $0\leq b_{m,j}\leq\lambda^j$ and $b_{m,j}\to\lambda^j$; this establishes attainability, not an optimal asymptotic yield. \sourcecite{ref:ser18-1}{BESP03}, \sourcecite{ref:ser18-2}{CE12}
\end{itemize}

\subsection*{References}

\begin{enumerate}[label={},leftmargin=4.8em,labelsep=0.5em]
  \footnotesize
  \item[\textup{[BESP03]}]\label{ref:ser18-1}
  D. E. Browne, J. Eisert, S. Scheel, and M. B. Plenio, "Driving Non-Gaussian to Gaussian States with Linear Optics," \emph{Physical Review A} \textbf{67}, 062320 (2003). \href{https://doi.org/10.1103/PhysRevA.67.062320}{doi:10.1103/PhysRevA.67.062320}; \href{https://arxiv.org/abs/quant-ph/0211173}{arXiv:quant-ph/0211173}.

  \item[\textup{[CE12]}]\label{ref:ser18-2}
  E. T. Campbell and J. Eisert, "Gaussification and Entanglement Distillation of Continuous Variable Systems: A Unifying Picture," \emph{Physical Review Letters} \textbf{108}, 020501 (2012). \href{https://doi.org/10.1103/PhysRevLett.108.020501}{doi:10.1103/PhysRevLett.108.020501}; \href{https://arxiv.org/abs/1107.1406}{arXiv:1107.1406}.
\end{enumerate}

\subsection*{Comment}

The optimal probability for arbitrary non-Gaussian inputs remains unresolved. The displayed convergent protocol has exponentially decreasing success probability when every branch of a fixed input tree must succeed. Convergence alone does not establish a positive asymptotic yield.

\subsection*{Field}

Quantum Resource Theory

\subsection*{Topic}

Gaussian quantum information; Entanglement distillation; One-shot and finite-blocklength bounds

\subsection*{ID}

\texttt{op\_6f6de1a285416bb2}
